% Packages are similar to Python dependencies: each adds a focused capability.
\usepackage[a4paper,top=0.42in,bottom=0.42in,left=0.48in,right=0.48in]{geometry}
\usepackage[T1]{fontenc}
\usepackage{lmodern}
\usepackage{microtype}
\usepackage{xcolor}
\usepackage{hyperref}
\usepackage{enumitem}
\usepackage{tabularx}
\usepackage{array}
\usepackage{titlesec}
\usepackage{needspace}

\definecolor{accent}{HTML}{145DA0}
\definecolor{rulegray}{HTML}{7A7A7A}

\hypersetup{
  colorlinks=true,
  urlcolor=accent,
  linkcolor=accent,
  pdfauthor={Gabriel Torres},
  pdftitle={Gabriel Torres - Resume}
}

\pagestyle{empty}
\setlength{\parindent}{0pt}
\setlength{\parskip}{0pt}
\setlength{\tabcolsep}{0pt}
\urlstyle{same}

\titleformat{\section}
  {\large\bfseries}
  {}{0pt}{}
  [\vspace{-0.35em}\color{rulegray}\titlerule]
\titlespacing*{\section}{0pt}{0.55em}{0.3em}

\setlist[itemize]{
  leftmargin=1.35em,
  itemsep=0.03em,
  parsep=0pt,
  topsep=0.12em,
  partopsep=0pt
}

% Commands are formatter functions. Their arguments are the resume's data.
\newcommand{\resumeheader}[5]{%
  \begin{tabularx}{\textwidth}{@{}p{0.30\textwidth}>{\centering\arraybackslash}X>{\raggedleft\arraybackslash}p{0.30\textwidth}@{}}
    #1 & {\LARGE\bfseries #3} & #4 \\
    #2 & \hfill & #5
  \end{tabularx}
}

\newcommand{\experience}[4]{%
  \Needspace{4\baselineskip}%
  \vspace{0.35em}%
  \begin{tabularx}{\textwidth}{@{}>{\raggedright\arraybackslash}p{0.38\textwidth}@{\hspace{0.8em}}>{\centering\arraybackslash}X@{\hspace{0.8em}}>{\raggedleft\arraybackslash}p{0.16\textwidth}@{}}
    #1 & \textbf{#2} & #3 \\
       & #4 &
  \end{tabularx}%
  \vspace{-0.2em}%
}

\newcommand{\project}[2]{\item \textbf{\textcolor{accent}{#1:}} #2}
\newcommand{\skillrow}[2]{\item \textbf{#1:} #2}
\newcommand{\educationitem}[2]{\item \textbf{#1}, #2}
